\begin{textsample}{POS Dim 1 – vem\_ed\_20 – Score 3.00 – vem\_ed\_20/t098.txt}  \label{ex:f1_pos_112}
Línguas e Aprendizagem Global Esse \textbf{foi} o tema do VI CBTecLE, Congresso Brasileiro de Línguas na Formação Técnica e Tecnológica, realizado presencialmente na Fatec Guaratinguetá entre 13 e 15 de setembro de 2023.
A equipe de PCIs / Cesu apoiou a organização do evento.
A abertura, no dia 14, \textbf{contou} com a participação do coordenador técnico da Cesu, Rafael Ferreira Alves; do diretor da Fatec Guaratinguetá, André Amarante; do coordenador dos PCIs / Cesu, Osvaldo Succi Junior; da coordenadora do Eixo de Línguas e Projetos Internacionais, Mariane Teixeira e com a representante da ARInter, Renata Rezende.
O coordenador técnico da Cesu ressaltou a \textbf{importância} das línguas para a empregabilidade e \textbf{destacou} o crescimento dos PCIs.
Osvaldo Succi Junior \textbf{destacou} os \textbf{números} dos PCIs: cerca de 60 projetos por semestre, envolvendo mais de 1.100 alunos das Fatecs em colaborações em inglês, espanhol e português.
Mariane Teixeira retomou a história do CBTecLE,"único no Brasil voltado para ensino de línguas para formação técnica e tecnológica"; agradeceu a organização da Fatec Guaratinguetá e frisou \textbf{que} línguas e internacionalização estão sempre juntas.
Colaboradoras internacionais de PCIs com Fatecs, Anoush Ayunts (Yerevan State University, Armênia) e Blanca Betancourt (Uniminuto, Colômbia) apresentaram a mesa redonda"Formação de Professores para o Ensino de Línguas"com Elza Maria Melo Ribeiro (IFRJ).
A mediação \textbf{foi} de Edilene Gasparini (CPS / Fatec São José do Rio Preto).
Anoush Ayunts fez uma breve apresentação sobre seu país, a Armênia, declarou-se apaixonada pelo Brasil e explicou sobre o PCI realizado entre Yerevan State University e as Fatecs Bragança Paulista, Garça, Guaratinguetá, Jales, Mogi das Cruzes, Pompeia e São José do Rio Preto.
Anoush e Edilene coordenaram as atividades do projeto, \textbf{que} adotou a metodologia Problem Based Learning e o uso da ferramenta Webquest para abordagem de problemas reais, como a fome, a seca, a dependência de tecnologias.

% matched lemmas: contar, destacar, importância, número, que, ser
\end{textsample}
